\chapter{Twierdzenie Poincarégo w wymiarze trzy: droga przez potok Ricciego}
\label{ch:poincare-trzy}
\index{twierdzenie Poincarego@Twierdzenie Poincarégo!wymiar trzy}

W Rozdziale~\ref{ch:poincare-2-high} poznaliśmy przypadek
dwuwymiarowy i wysokowymiarowy. Wymiar trzy wymaga innych
narzędzi. Zaczniemy od dowolnej metryki na badanej
rozmaitości, będziemy ją zmieniać potokiem Ricciego
z Rozdziału~\ref{ch:potok-ricciego}, a w miejscach powstawania
osobliwości wykonamy kontrolowane cięcia. Na końcu
odczytamy typ rozmaitości z tego, co zniknęło podczas
ewolucji.

Ten rozdział objaśnia \emph{drogę dowodu}, lecz nie jest
pełnym dowodem twierdzeń Perelmana. Własne rachunki i krótkie
dowody topologiczne podajemy w całości. Twierdzenia o braku
zapadania, otoczeniach kanonicznych, istnieniu potoku
z chirurgią i wygasaniu w skończonym czasie oznaczamy
jako \emph{wejścia zewnętrzne}. Ich dowody wymagają analizy
potoku, oszacowań geometrycznych i teorii powierzchni
minimalnych wykraczających poza ten skrypt. Pełne opracowanie
tych kroków zawiera monografia Morgana i Tiana wskazana
na końcu rozdziału.

\begin{twierdzenie}[Poincaré w wymiarze trzy]
\label{thm:poincare-trzy}
Każda gładka, zamknięta, spójna i jednospójna
$3$-rozmaitość jest dyfeomorficzna z $S^3$.
\end{twierdzenie}

Słowo \emph{zamknięta} oznacza zwarta i bez brzegu.
Jednospójność wyklucza przeszkody widoczne w grupie
podstawowej, ale sama nie pozwala od razu zbudować
dyfeomorfizmu. Teoria z Rozdziału~\ref{ch:chirurgia-geometryczna-21}
również nie rozwiązuje tego zagadnienia automatycznie:
jej argumenty wysokowymiarowe używają ruchów rozłączania,
których nie można tu po prostu przenieść. Potok Ricciego
wnosi dodatkową informację geometryczną o miejscach,
w których należy ciąć.

\section{Dlaczego ewolucja metryki pomaga?}
\label{sec:poincare3-pierwsze-rachunki}

Przypomnijmy równanie nienormalizowanego potoku Ricciego:
\begin{equation}\label{eq:poincare3-ricci}
 \partial_tg=-2\operatorname{Ric}(g).
\end{equation}
Twierdzenie~\ref{thm:ricci-short-time} zapewnia jego
gładkie rozwiązanie tylko przez krótki czas. Kształt
przyszłych osobliwości zależy od geometrii całej
rozmaitości. Wymiar trzy jest jednak szczególny: tensor
Ricciego zawiera w nim całą informację o krzywiźnie.

\begin{stwierdzenie}[Krzywizna przekroju w wymiarze trzy]
\label{prop:poincare3-ricci-sekcje}
Niech $(e_1,e_2,e_3)$ będzie bazą ortonormalną
$T_pM^3$, a $K_{ij}$ krzywizną przekroju rozpiętego
przez $e_i,e_j$. Wtedy
\begin{equation}\label{eq:poincare3-ricci-sekcje}
 K_{12}=\frac{\operatorname{Ric}(e_1,e_1)
              +\operatorname{Ric}(e_2,e_2)
              -\operatorname{Ric}(e_3,e_3)}{2}.
\end{equation}
Pozostałe dwa wzory otrzymujemy przez zamianę indeksów.
\end{stwierdzenie}
\begin{proof}
Z definicji Ricciego jako śladu krzywizny mamy
\[
 \operatorname{Ric}(e_1,e_1)=K_{12}+K_{13},\quad
 \operatorname{Ric}(e_2,e_2)=K_{12}+K_{23},\quad
 \operatorname{Ric}(e_3,e_3)=K_{13}+K_{23}.
\]
Dodanie dwóch pierwszych równości i odjęcie trzeciej
daje $2K_{12}$, czyli~\eqref{eq:poincare3-ricci-sekcje}.
\end{proof}

Wybór dowolnej bazy ortonormalnej pozwala odzyskać
krzywiznę dowolnej płaszczyzny z Ricciego. Rachunek
nie oznacza, że krzywizna ma zawsze jeden znak:
trzy wartości własne Ricciego mogą być różne.
Pokazuje natomiast, czemu kontrola Ricciego jest
szczególnie silna właśnie w tym wymiarze.

\begin{przyklad}[Kurcząca się sfera]
\label{prz:poincare3-kurczaca-sfera}
Dla $g_0=a_0^2g_{S^3}$, gdzie $g_{S^3}$ ma
krzywiznę sekcyjną jeden, wzór~\eqref{eq:ricci-sphere}
daje
\[
 g(t)=(a_0^2-4t)g_{S^3},\qquad
 0\leq t<\frac{a_0^2}{4}.
\]
Promień maleje do zera; po przeskalowaniu metryki
przez czynnik $(a_0^2-4t)^{-1}$ wciąż widzimy
okrągłą sferę. Gładkie rozwiązanie kończy się
w skończonym czasie, chociaż topologia sfery jest
całkowicie rozpoznawalna.
\end{przyklad}

\begin{przyklad}[Kurczący się walec]
\label{prz:poincare3-walec}
Na $S^2\times\R$ weźmy
$g_0=a_0^2g_{S^2}+\dd z^2$. Iloczyn ma
\[
 \operatorname{Ric}(a^2g_{S^2}+\dd z^2)
 =g_{S^2}+0\cdot\dd z^2.
\]
Stąd bezpośrednie wstawienie do~\eqref{eq:poincare3-ricci}
daje
\begin{equation}\label{eq:poincare3-walec}
 g(t)=(a_0^2-2t)g_{S^2}+\dd z^2,
 \qquad 0\leq t<\frac{a_0^2}{2}.
\end{equation}
Przekroje $S^2$ kurczą się, a kierunek osi pozostaje
niezmieniony. Krzywizna skalarna równa się
$2/(a_0^2-2t)$ i dąży do nieskończoności.
Po zastąpieniu osi $\R$ przez okrąg dostajemy
analogiczny przykład zamkniętej rozmaitości
$S^2\times S^1$.
\end{przyklad}

Model walca wyjaśnia słowo \emph{szyjka}:
w obszarze przypominającym $S^2\times(-L,L)$
przekroje sferyczne stają się bardzo małe.
Nie każda osobliwość jest dokładnym walcem;
oprócz szyjek pojawiają się na przykład
zakończenia podobne do walca zamkniętego
z jednej strony. Ich rozpoznanie jest jednym
z głównych wyników użytych poniżej.

Równanie skalarne~\eqref{eq:ricci-scalar-evolution}
pozwala wykonać jeszcze jeden użyteczny rachunek.

\begin{stwierdzenie}[Dolne oszacowanie skalara]
\label{prop:poincare3-skalar-dol}
Niech $g(t)$ będzie gładkim potokiem Ricciego
na zamkniętej $3$-rozmaitości i niech
$\operatorname{Scal}_{g(0)}\geq-C$, gdzie $C>0$.
W całym przedziale istnienia zachodzi
\begin{equation}\label{eq:poincare3-skalar-dol}
 \operatorname{Scal}_{g(t)}
 \geq-\frac{3}{2(t+3/(2C))}.
\end{equation}
\end{stwierdzenie}
\begin{proof}
W ortonormalnej bazie $\operatorname{Ric}$ ma trzy
rzeczywiste wartości własne $\lambda_i$. Nierówność
Cauchy'ego--Schwarza daje
\[
 |\operatorname{Ric}|^2=\sum_i\lambda_i^2
 \geq\frac13\Bigl(\sum_i\lambda_i\Bigr)^2
 =\frac13\operatorname{Scal}^2.
\]
Z~\eqref{eq:ricci-scalar-evolution} wynika więc
$\partial_t\operatorname{Scal}
 \geq\Delta_g\operatorname{Scal}
 +(2/3)\operatorname{Scal}^2$.
Funkcja stała w przestrzeni
$u(t)=-3/(2(t+3/(2C)))$ spełnia
$u'=(2/3)u^2$ i $u(0)=-C$.
Zasada maksimum na zwartej rozmaitości,
zastosowana do dolnego rozwiązania $u$,
daje $\operatorname{Scal}\geq u$.
\end{proof}

Oszacowanie mówi, że duża ujemna krzywizna
skalarna nie może pojawić się nagle w gładkim
potoku. Nie ogranicza jednak krzywizny od góry:
na kurczącej się sferze i walcu skalar rośnie
do nieskończoności. Kontrola potoku przez
osobliwości wymaga innych narzędzi.

\section{Powiększanie osobliwości i brak zapadania}
\label{sec:poincare3-osobliwosci}
\index{potok Ricciego!brak zapadania}

Załóżmy, że w ciągu punktów i chwil $(x_i,t_i)$
wartość $Q_i=|\operatorname{Rm}|(x_i,t_i)$
dąży do nieskończoności. Aby obejrzeć geometrię
w skali krzywizny, wprowadzamy nowe metryki
\begin{equation}\label{eq:poincare3-rescale}
 g_i(s)=Q_i g(t_i+s/Q_i).
\end{equation}
W punkcie bazowym ich krzywizna ma rozmiar jeden.
Jest to \emph{powiększenie paraboliczne}:
odległości rosną $Q_i^{1/2}$ razy, a jednostka
czasu $Q_i$ razy. Równość~\eqref{eq:poincare3-ricci}
pozostaje prawdziwa, ponieważ Ricci jako tensor
typu $(0,2)$ nie zmienia się przy stałej zmianie
skali, a pochodna lewej strony daje
$\partial_sg_i=\partial_tg$.

Ograniczona krzywizna w tej nowej skali nie
gwarantuje jeszcze gładkiej granicy. Mogłaby
zapaść się objętość kul. Przykładowo iloczyn
$\R^2\times S^1_\varepsilon$ z płaską metryką
ma krzywiznę zero, ale kula o promieniu jeden
ma objętość rzędu $\varepsilon$. Gdy
$\varepsilon\to0$, nie zbiega ona do
trójwymiarowej kuli euklidesowej.

\begin{definicja}[Brak zapadania na danej skali]
\label{def:poincare3-noncollapsing}
Potok trójwymiarowy nazywamy
$\kappa$-\emph{niezapadającym się} na skalach
co najwyżej $\rho$, jeśli dla każdego
$0<r\leq\rho$ i każdego punktu $(x,t)$,
dla którego wsteczna kula czasoprzestrzenna
\[
 P(x,t,r,-r^2)=
 \{(y,s):t-r^2\leq s\leq t,
        \ d_{g(t)}(x,y)<r\}
\]
leży w gładkim obszarze potoku i spełnia
$|\operatorname{Rm}|\leq r^{-2}$,
zachodzi
\begin{equation}\label{eq:poincare3-noncollapsing}
 \operatorname{Vol}_{g(t)}B_{g(t)}(x,r)
 \geq\kappa r^3.
\end{equation}
Definicja opisuje lokalną objętość względem
promienia. Przy potoku z chirurgią jej użycie
wymaga doprecyzowania, które kule nie przecinają
czasów cięcia.
\end{definicja}

\begin{twierdzenie}[Perelman, twierdzenie 4.1:
 brak lokalnego zapadania; wejście zewnętrzne]
\label{thm:poincare3-perelman-noncollapse}
Niech $g(t)$ będzie gładkim potokiem Ricciego
na zamkniętej $3$-rozmaitości $M$, określonym
na $[0,T)$ z $T<\infty$. Nie istnieją ciągi
$t_i\to T$, punktów $x_i\in M$ i promieni
$r_i>0$, dla których jednocześnie
$r_i^2/t_i$ pozostaje ograniczone,
$|\operatorname{Rm}|\leq r_i^{-2}$
na $B_{g(t_i)}(x_i,r_i)$ oraz
\[
 r_i^{-3}\operatorname{Vol}_{g(t_i)}
 B_{g(t_i)}(x_i,r_i)\longrightarrow0.
\]
\end{twierdzenie}

Jest to dokładnie postać z $\S\,4.1$ pracy
\href{https://arxiv.org/abs/math/0211159}
{G.\ Perelmana, \emph{The entropy formula for the
Ricci flow and its geometric applications} (2002)}.
Wersja~\eqref{eq:poincare3-noncollapsing}
używa ograniczenia krzywizny w całej kuli
czasoprzestrzennej. W twierdzeniu dla ciągów
$t_i\to T$ wystarcza założenie o krzywiźnie
tylko w chwili $t_i$; to odrębna, mocniejsza
w tym punkcie postać braku zapadania.
Kontrola potrzebna przy \emph{chirurgii} jest
odrębnym krokiem, rozwiniętym w
\href{https://arxiv.org/abs/math/0303109}
{drugim preprincie Perelmana (2003)} oraz
w \href{https://www.claymath.org/library/monographs/cmim03.pdf}
{Morgan--Tian, rozdz.\ 8 i 16}.

Skąd bierze się takie twierdzenie? Perelman
wprowadził wielkości monotoniczne dla potoku.
Jedną z nich jest entropia
\[
 \mathcal W(g,f,\tau)=
 \int_M\left[\tau\bigl(|\nabla f|^2
                  +\operatorname{Scal}\bigr)+f-3\right]
 (4\pi\tau)^{-3/2}e^{-f}\,\dd\mu_g,
 \quad
 \int_M(4\pi\tau)^{-3/2}e^{-f}\,\dd\mu_g=1.
\]
Przy odpowiednim równaniu sprzężonym dla $f$
i $\tau'=-1$ jej pochodna ma postać
\begin{equation}\label{eq:poincare3-entropy}
 \frac{\dd}{\dd t}\mathcal W
 =2\tau\int_M
 \left|\operatorname{Ric}+\nabla^2f
              -\frac{g}{2\tau}\right|^2
 (4\pi\tau)^{-3/2}e^{-f}\,\dd\mu_g\geq0.
\end{equation}
Formuła i wykorzystanie jej do oszacowania objętości
są tu \emph{wejściem zewnętrznym}. Jej sens można
jednak odczytać bez dowodu: prawa strona nie jest
ujemna, więc pewna ilość informacji geometrycznej
nie ginie podczas ewolucji. Sam wzór nie kontroluje
automatycznie skoków metryki przy chirurgii;
ten etap wymaga osobnych oszacowań.

\section{Szyjki, czapy i otoczenia kanoniczne}
\label{sec:poincare3-kanoniczne}

\begin{definicja}[Szyjka $\varepsilon$]
\label{def:poincare3-neck}
Niech $\varepsilon>0$ będzie małe. Obszar $N$
w trójwymiarowej rozmaitości jest
\emph{szyjką $\varepsilon$}, jeżeli istnieje
dyfeomorfizm
\[
 \Phi:S^2\times(-\varepsilon^{-1},
                       \varepsilon^{-1})\longrightarrow N
\]
oraz dodatni czynnik skali $Q$, dla których
$Q\Phi^*g$ jest $\varepsilon$-bliskie metryce
okrągłego walca w odpowiednio wysokiej normie
$C^k$, na przykład dla
$k=\lfloor\varepsilon^{-1}\rfloor$.
Sfera $\Phi(S^2\times\{0\})$ jest
\emph{przekrojem centralnym}.
\end{definicja}

Definicja zawiera zarówno informację geometryczną,
jak i topologiczną. Bliskość metryk mówi, że
przekroje są prawie okrągłe; dyfeomorfizm mówi,
że przekrój jest rzeczywiście zanurzoną sferą,
a nie tylko powierzchnią wyglądającą podobnie
w jednym obrazie. Model~\eqref{eq:poincare3-walec}
jest szyjką dokładną po dobraniu skali.

\begin{definicja}[Czapa; opis topologiczny]
\label{def:poincare3-cap}
\emph{Czapą} nazywamy obszar z jednym końcem
szyjkowym, którego zwarty rdzeń jest dyfeomorficzny
z $D^3$ albo z $\mathbb{RP}^3\setminus
\operatorname{int}D^3$. Granica rdzenia jest
$S^2$. Dla czapy kanonicznej wymagane są ponadto
oszacowania krzywizny, średnicy i pochodnych
metryki na skali szyjki; sama postać topologiczna
nie wystarcza do bezpiecznej chirurgii.
\end{definicja}

Przykładem topologicznym jest półnieskończony
walec $S^2\times[0,\infty)$ zamknięty dyskiem
$D^3$ przy lewym końcu i wygładzony przy
sklejeniu. Geometria czapy używanej w potoku
jest dobierana starannie, by po sklejeniu
można było wznowić równanie~\eqref{eq:poincare3-ricci}.

Twierdzenie o \emph{otoczeniach kanonicznych}
rozpoznaje punkty o dużej krzywiźnie w potoku
z chirurgią. Po przeskalowaniu mogą one należeć
do szyjki, czapy albo zwartej składowej
o kontrolowanej dodatniej krzywiźnie.
Jest to \emph{opis jakościowy} wejścia
Perelmana, nie jego pełne sformułowanie
z parametrami $r,\delta,\varepsilon$.
W dalszym rozumowaniu korzystamy jedynie
z precyzyjnych konsekwencji dla istnienia
potoku i odtworzenia topologii, które
podajemy jako osobne twierdzenia.
Pełne hipotezy i oszacowania znajdują się
w \href{https://arxiv.org/abs/math/0211159}
{pierwszym} i
\href{https://arxiv.org/abs/math/0303109}
{drugim preprincie Perelmana}, a szczegółowy
dowód w
\href{https://www.claymath.org/library/monographs/cmim03.pdf}
{Morgan--Tian, rozdz.\ 9--11 i 17}.

Warto zobaczyć zależność między twierdzeniami.
Powiększenia~\eqref{eq:poincare3-rescale} mają
krzywiznę unormowaną, lecz mogłyby się zapaść.
Brak zapadania umożliwia uzyskanie geometrycznych
granic. Granicami są stare rozwiązania potoku,
zwane $\kappa$-rozwiązaniami, z nieujemną krzywizną
i dodatnimi dolnymi oszacowaniami objętości.
Ich analiza prowadzi do rozpoznania szyjek i czap.
Każda strzałka w tym łańcuchu wymaga osobnego
głębokiego twierdzenia; sama obserwacja o skalowaniu
nie dowodzi klasyfikacji granic.

\section{Cięcie szyjki i kontrola topologii}
\label{sec:poincare3-chirurgia}
\index{potok Ricciego!chirurgia}

Wybierzmy centralną sferę $S^2$ w bardzo cienkiej,
dobrze kontrolowanej szyjce. Przecinamy
rozmaitość wzdłuż niej, usuwamy obszary
o nadmiernej krzywiźnie i przyklejamy
do nowych sfer brzegowych standardowe
metryczne czapy. Jest to \emph{chirurgia
potoku Ricciego}. Od tej chwili rozwiązujemy
zwykłe równanie Ricciego z nową metryką,
aż nadejdzie kolejny czas chirurgii.

Operacja topologiczna przy pojedynczej sferze
jest szczególnym przypadkiem chirurgii
z Rozdziału~\ref{ch:chirurgia-geometryczna-21}:
wycinamy $S^2\times(-1,1)$ i do dwóch
powstałych brzegów przyklejamy $D^3$.
Różnica leży w warunku geometrycznym:
cięcie musi nastąpić w szyjce o kontrolowanej
skali, a nowe czapy muszą pozwalać na dalszy potok.

\begin{lemat}[Topologia cięcia wzdłuż sfery]
\label{lem:poincare3-sphere-cut}
Niech $M$ będzie zamkniętą, spójną,
zorientowaną $3$-rozmaitością, a
$S\subset M$ zanurzoną, dwustronną sferą.
Niech $N$ powstanie przez rozcięcie $M$
wzdłuż $S$ i zamknięcie nowych brzegów
dwoma dyskami $D^3$.
Jeżeli $S$ rozdziela $M$, to $N=N_1\sqcup N_2$
i $M\cong N_1\#N_2$.
Jeżeli $S$ nie rozdziela $M$, to $N$ jest spójna
i $M\cong N\#(S^2\times S^1)$.
Wszystkie równoważności są dyfeomorfizmami.
\end{lemat}
\begin{proof}
Po wycięciu małego iloczynowego otoczenia
$S^2\times(-1,1)$ otrzymujemy rozmaitość
z dwiema sferami brzegowymi. W przypadku
rozdzielającym leżą one w różnych składowych.
Zamknijmy każdą brzegową sferę dyskiem.
Odwracając operację, usuwamy wnętrza tych
dwóch dysków i sklejamy otrzymane brzegi
walcem $S^2\times[0,1]$. Jest to konstrukcja
sumy spójnej $N_1\#N_2$.

Gdy dopełnienie $S$ jest spójne, obie sfery
brzegowe należą do tej samej składowej.
Po ich zamknięciu dostajemy $N$.
Odtworzenie $M$ polega na usunięciu wnętrz
dwóch dysków z $N$ i połączeniu ich sfer
brzegowych walcem. Ta operacja dokłada
jeden $1$-uchwyt do $N$ i daje
$N\#(S^2\times S^1)$. Orientacja $M$
wyklucza skręconą wiązkę $S^2$ nad $S^1$.
Równoważność można też sprawdzić przez
wykonanie tego samego cięcia na standardowym
$S^2\times S^1$ wzdłuż włókna $S^2$.
\end{proof}

\begin{uwaga}[Co mogą zrobić rzeczywiste chirurgie]
\label{rem:poincare3-discarded}
Powyższy lemat opisuje pojedyncze cięcie,
ale w potoku usuwa się również całe składowe,
których geometria została rozpoznana. W przypadku
orientowalnym mogą to być sferyczne formy
przestrzenne $S^3/\Gamma$, $S^2\times S^1$
oraz $\mathbb{RP}^3\#\mathbb{RP}^3$.
Ostatni przypadek jest już sumą spójną
dwóch sferycznych form. Czapę o rdzeniu
$\mathbb{RP}^3\setminus\operatorname{int}D^3$
można również zastąpić dyskiem; topologicznie
odpowiada to odłączeniu czynnika
$\mathbb{RP}^3$.
Dokładna lista i jej uzasadnienie są częścią
zewnętrznej kontroli otoczeń kanonicznych.
\end{uwaga}

\begin{twierdzenie}[Perelman: kontrolowany potok
 z chirurgią; wejście zewnętrzne]
\label{thm:poincare3-surgery-flow}
Dla każdej gładkiej metryki początkowej
na zamkniętej, zorientowanej $3$-rozmaitości
można skonstruować potok Ricciego z chirurgią
na wszystkich czasach $t\geq0$.
Na przedziałach między chirurgiami spełnia
on~\eqref{eq:poincare3-ricci}; w każdym
ograniczonym przedziale czasu zachodzi tylko
skończenie wiele chirurgii. Przejście przez
czas chirurgii zmienia topologię przez
rozkład sumy spójnej i usunięcie rozpoznanych
składowych. W przypadku zorientowanym
usuwane składowe są dyfeomorficzne do
$S^2\times S^1$, $\mathbb{RP}^3\#\mathbb{RP}^3$
albo sferycznej formy przestrzennej.
\end{twierdzenie}

W przywołanej monografii
\href{https://www.claymath.org/library/monographs/cmim03.pdf}
{Morgan--Tian, tw.\ 0.3 i 15.9}
twierdzenie o istnieniu jest sformułowane
przy wykluczeniu zanurzonej obustronnej
$\mathbb{RP}^2$. Każda zorientowana
$3$-rozmaitość spełnia ten warunek:
dwustronna podrozmaitość kowymiaru jeden
w zorientowanej rozmaitości odziedziczyłaby
orientację, a $\mathbb{RP}^2$ jest
nieorientowalna. Otrzymujemy więc powyższą
wersję potrzebną w tym rozdziale.
Pierwotnym źródłem konstrukcji jest
\href{https://arxiv.org/abs/math/0303109}
{Perelman, \emph{Ricci flow with surgery on
three-manifolds} (2003)}.
W szczegółowej konstrukcji metrykę początkową
normalizuje się stałym przeskalowaniem,
a progi cięcia dobiera tak, by zachować
oszacowania krzywizny i niezapadania.

Sam opis „tnij szyjkę i przyklej czapę”
nie dowodzi twierdzenia o istnieniu.
Trzeba zapewnić, że progi cięcia można
wybierać w kolejnych etapach, czapy
zachowują wymagane oszacowania,
nowe osobliwości znów mają otoczenia
kanoniczne i czasy cięć nie gromadzą się
w skończonym przedziale. To właśnie
stanowi znaczną część dowodu Perelmana.

\section{Dlaczego potok musi wygasnąć?}
\label{sec:poincare3-wygasanie}
\index{potok Ricciego!wygaśnięcie}

Potok z chirurgią może być określony dla
wszystkich $t\geq0$, chociaż po pewnym
czasie przekrój $M_t$ staje się pusty.
To nazywamy \emph{wygaśnięciem w skończonym
czasie}. Nie znaczy to, że każda metryka
na każdej $3$-rozmaitości znika:
płaski torus z Przykładu~\ref{prz:ricci-torus}
jest rozwiązaniem stacjonarnym.

\begin{twierdzenie}[Perelman: wygasanie;
 wejście zewnętrzne]
\label{thm:poincare3-extinction}
Niech $M$ będzie zamkniętą, spójną,
zorientowaną $3$-rozmaitością, dla której
$\pi_1(M)$ jest wolnym iloczynem skończenie
wielu grup skończonych i kopii $\Z$.
Dla każdej gładkiej metryki początkowej
kontrolowany potok Ricciego z chirurgią,
skonstruowany jak w
Twierdzeniu~\ref{thm:poincare3-surgery-flow},
ma czas $T<\infty$, po którym $M_t=\varnothing$
dla każdego $t\geq T$.
\end{twierdzenie}

Jest to wersja zastosowana w
\href{https://www.claymath.org/library/monographs/cmim03.pdf}
{Morgan--Tian, tw.\ 18.1}; dowód zajmuje
tam rozdziały 18--19. W
\href{https://arxiv.org/abs/math/0307245}
{trzecim preprincie Perelmana,
\emph{Finite extinction time for the solutions
to the Ricci flow on certain three-manifolds}
(2003)} hipoteza została wyrażona przez brak
czynników asferycznych w rozkładzie pierwszym.
W tym rozdziale korzystamy z wersji grupowej,
bo dla $\pi_1(M)=1$ jej założenie widać od razu.

Idea wygasania ma mierzalny ślad w rachunku
powierzchni. Poniższa nierówność dotyczy
\emph{gładkiego odcinka} potoku, bez chirurgii.

\begin{stwierdzenie}[Zmiana pola minimalnej sfery]
\label{prop:poincare3-minimal-area}
Niech $F:S^2\to(M^3,g(t_0))$ będzie gładkim
minimalnym zanurzeniem w chwili $t_0$.
Trzymając $F$ nieruchomo i zmieniając tylko
metrykę potokiem Ricciego, otrzymujemy
\begin{equation}\label{eq:poincare3-minimal-area}
 \left.\frac{\dd}{\dd t}\right|_{t=t_0}
 \operatorname{Area}_{g(t)}(F)
 =-4\pi-\frac12\int_{S^2}
   \bigl(\operatorname{Scal}+|A|^2\bigr)\,\dd\mu_F
 \leq-4\pi-\frac12R_{\min}(t_0)
                 \operatorname{Area}_{g(t_0)}(F),
\end{equation}
gdzie $A$ jest drugą formą fundamentalną
zanurzenia, a
$R_{\min}(t_0)=\min_M\operatorname{Scal}_{g(t_0)}$.
\end{stwierdzenie}
\begin{proof}
Pochodna miary powierzchni indukowanej
przez $g(t)$ wynosi
$\frac12\operatorname{tr}_{TF}(\partial_tg)$
razy ta miara. Z~\eqref{eq:poincare3-ricci}
dostajemy więc
\[
 \frac{\dd}{\dd t}\operatorname{Area}_{g(t)}(F)
 =-\int_{S^2}\operatorname{tr}_{TF}
                    \operatorname{Ric}\,\dd\mu_F.
\]
Wybierzmy lokalną bazę ortonormalną
$(e_1,e_2,\nu)$, gdzie $\nu$ jest normalną.
Z definicji skalara i Ricciego
\[
 \operatorname{tr}_{TF}\operatorname{Ric}
 =\operatorname{Scal}-\operatorname{Ric}(\nu,\nu)
 =\tfrac12\operatorname{Scal}+K_M(e_1,e_2).
\]
Równanie Gaussa dla zanurzonej powierzchni daje
$K_F=K_M(e_1,e_2)+\det A$.
Minimalność oznacza, że krzywizny główne
są przeciwne, stąd $\det A=-|A|^2/2$.
Zatem
\[
 \operatorname{tr}_{TF}\operatorname{Ric}
 =\tfrac12\operatorname{Scal}+K_F
                   +\tfrac12|A|^2.
\]
Całkujemy i używamy twierdzenia Gaussa--Bonneta
$\int_{S^2}K_F\,\dd\mu_F=4\pi$.
Nierówność wynika z $|A|^2\geq0$ oraz
$\operatorname{Scal}\geq R_{\min}$.
\end{proof}

Jeżeli istniałaby przez dowolnie długi czas
rodzina istotnych sfer minimalnych, której
pole spełnia odpowiednią wersję
\eqref{eq:poincare3-minimal-area}, to
oszacowanie~\eqref{eq:poincare3-skalar-dol}
prowadziłoby do nierówności typu
\[
 W'(t)\leq-4\pi+
 \frac{3}{4(t+c)}W(t),\qquad c>0.
\]
Po pomnożeniu przez $(t+c)^{-3/4}$
pochodna prawego iloczynu jest nie większa
niż $-4\pi(t+c)^{-3/4}$. Całka z tej
ujemnej funkcji ma moduł rosnący bez granic,
więc w końcu wymusiłaby $W(t)<0$,
niemożliwe dla pola lub energii.

To obliczenie pokazuje mechanizm, ale nie jest
dowodem Twierdzenia~\ref{thm:poincare3-extinction}.
Trzeba skonstruować właściwą wielkość min--max,
uzyskać jej oszacowanie pochodnej mimo możliwego
powstawania bąbli w granicach powierzchni
minimalnych oraz wykazać, że chirurgia nie
psuje potrzebnej nierówności. Perelman używa
rodzin pętli i minimalnych dysków je wypełniających;
Morgan i Tian rozpisują ten argument. Istnieje
także podejście przez rodziny odwzorowań $S^2$,
opisane w tej monografii jako wariant
Colding--Minicozziego.

\begin{lemat}[Która grupa homotopii jest istotna?]
\label{lem:poincare3-pi3}
Jeżeli $M$ jest zamkniętą i jednospójną
$3$-rozmaitością, to
$\pi_2(M)=0$ oraz $\pi_3(M)\cong\Z$.
\end{lemat}
\begin{proof}
Jednospójność daje orientowalność: gdyby
wiązka orientacji była nietrywialna, powstałoby
niebanalne dwukrotne nakrycie. Mamy $H_1(M;\Z)=0$,
więc z twierdzenia o współczynnikach
uniwersalnych $H^1(M;\Z)=0$.
Dualność Poincarégo
(Twierdzenie~\ref{thm:poincare-dualnosc-16})
daje $H_2(M;\Z)\cong H^1(M;\Z)=0$
oraz $H_3(M;\Z)\cong\Z$.
Ponieważ $M$ jest jednospójna,
Twierdzenie~\ref{thm:hurewicz-wyzszy}
w najniższym możliwym stopniu daje
$\pi_2(M)\cong H_2(M)=0$.
Stosujemy je ponownie, teraz w stopniu trzy,
i otrzymujemy $\pi_3(M)\cong H_3(M)=\Z$.
\end{proof}

Właśnie dlatego dowód dla jednospójnego $M$
nie może skończyć się na analizie pojedynczej
istotnej sfery $S^2\to M$: takiej klasy
homotopii nie ma. Niezerowy element $\pi_3(M)$
można przedstawić jako rodzinę sfer
$S^2\to M$ parametryzowaną przedziałem,
której końce są odwzorowaniami stałymi.
Jej parametr i sfera razem tworzą, po
odpowiednim sklejeniu, $S^3$.
Energia największej sfery w takiej rodzinie
prowadzi do analogicznej wielkości min--max.
U Perelmana odpowiada jej rodzina pętli
z minimalnymi dyskami rozpiętymi na tych
pętlach. Przejście od intuicji do nierówności
przez czasy chirurgii jest zewnętrzną częścią
twierdzenia o wygasaniu.

\section{Odtwarzanie topologii i końcowy wniosek}
\label{sec:poincare3-topologia}

Wygasanie mówi, że końcowy przekrój potoku jest
pusty. Aby dowiedzieć się, czym było $M_0$,
cofamy się przez skończenie wiele chirurgii
przed czasem wygasania. Z Lematu~\ref{lem:poincare3-sphere-cut}
wynika, że odwrócenie cięcia odtwarza sumę
spójną; Twierdzenie~\ref{thm:poincare3-surgery-flow}
ogranicza typy usuniętych składowych.

\begin{wniosek}[Odtworzenie po wygasaniu]
\label{thm:poincare3-reconstruction}
Niech $M$ będzie zamkniętą, spójną,
zorientowaną $3$-rozmaitością. Załóżmy, że
potok Ricciego z chirurgią, skonstruowany
jak w Twierdzeniu~\ref{thm:poincare3-surgery-flow},
ulega wygaśnięciu w skończonym czasie.
Wtedy $M$ jest dyfeomorficzna z sumą spójną
skończenie wielu sferycznych form
przestrzennych $S^3/\Gamma$ i kopii
$S^2\times S^1$. Grupa $\Gamma$ jest
skończoną grupą działającą swobodnie
i izometrycznie na okrągłej $S^3$.
\end{wniosek}
\begin{proof}
W przedziale od zera do chwili wygaśnięcia
jest skończenie wiele czasów chirurgii na mocy
Twierdzenia~\ref{thm:poincare3-surgery-flow}.
Między nimi gładki potok nie zmienia typu
dyfeomorficznego przekrojów. Cofając się
przez ostatnią chirurgię, odtwarzamy poprzedni
przekrój przez sumy spójne składowych
pozostałych po cięciu z usuniętymi składowymi.
Lemat~\ref{lem:poincare3-sphere-cut}
opisuje tę operację na przekrojach szyjek.
Usunięte składowe z listy w Twierdzeniu
~\ref{thm:poincare3-surgery-flow} są sferycznymi
formami, kopiami $S^2\times S^1$ albo
$\mathbb{RP}^3\#\mathbb{RP}^3$; ten ostatni
przypadek daje dwie sferyczne formy.
Powtarzając argument wstecz przez pozostałe
czasy chirurgii i zaczynając od pustego
przekroju końcowego, otrzymujemy podany
rozkład $M$.
\end{proof}

W pełnym twierdzeniu dla rozmaitości
nieorientowalnych pojawia się także skręcona
wiązka $S^2$ nad $S^1$. Wniosek powyżej
odpowiada
\href{https://www.claymath.org/library/monographs/cmim03.pdf}
{Morgan--Tian, wnioskowi 15.4},
po ograniczeniu do przypadku orientowalnego.
Zewnętrzna pozostaje klasyfikacja
\emph{wszystkich} usuwanych obszarów,
przyjęta w Twierdzeniu~\ref{thm:poincare3-surgery-flow}.
Rozpoznanie w
Uwadze~\ref{rem:poincare3-discarded}
wyjaśnia, czemu $\mathbb{RP}^3\#\mathbb{RP}^3$
nie potrzebuje osobnego rodzaju czynnika.

\begin{lemat}[Grupa podstawowa sumy spójnej]
\label{lem:poincare3-pi1-sum}
Dla spójnych, zamkniętych $3$-rozmaitości
$P$ i $Q$ zachodzi
\[
 \pi_1(P\#Q)\cong\pi_1(P)*\pi_1(Q).
\]
W szczególności
$\pi_1(S^2\times S^1)\cong\Z$ oraz
$\pi_1(S^3/\Gamma)\cong\Gamma$.
\end{lemat}
\begin{proof}
Suma spójna jest sklejeniem
$P\setminus\operatorname{int}D^3$ i
$Q\setminus\operatorname{int}D^3$ wzdłuż
sfery $S^2$. Usunięcie wnętrza dysku
nie zmienia grupy podstawowej:
odwracając operację, przyklejamy
jednospójny dysk wzdłuż jednospójnej
sfery. Teraz Twierdzenie~\ref{thm:van-kampen}
dla odpowiednich otwartych pogrubień
obu kawałków daje wolny iloczyn, bo ich
część wspólna odkształca się do $S^2$
i ma trywialną grupę podstawową.
Iloczyn $S^2\times S^1$ ma grupę $\Z$.
Ponieważ $S^3$ jest uniwersalnym nakryciem
$S^3/\Gamma$, jego grupa podstawowa
jest grupą przekształceń pokrycia $\Gamma$.
\end{proof}

\begin{proof}[Dowód Twierdzenia~\ref{thm:poincare-trzy}]
Niech $M$ spełnia hipotezy. Wybierzmy
dowolną gładką metrykę $g_0$.
Ponieważ $\pi_1(M)=1$, warunek grupowy
Twierdzenia~\ref{thm:poincare3-extinction}
jest spełniony: grupa trywialna jest
pustym wolnym iloczynem. Twierdzenia
~\ref{thm:poincare3-surgery-flow}
i~\ref{thm:poincare3-extinction} dają
kontrolowany potok z chirurgią, którego
przekrój staje się pusty w skończonym czasie.

Z Wniosku~\ref{thm:poincare3-reconstruction}
rozmaitość $M$ ma postać
\[
 M\cong\bigl(S^3/\Gamma_1\bigr)\#\cdots\#
 \bigl(S^3/\Gamma_k\bigr)
 \#\underbrace{(S^2\times S^1)\#\cdots\#
 (S^2\times S^1)}_{\ell\ \text{razy}}.
\]
Przyjmujemy tu, że $k+\ell\geq1$, ponieważ
$M$ jest niepusta i spójna; dla zerowej liczby
czynników w jednej z dwóch rodzin pomijamy
odpowiednią część zapisu.
Stosując Lemat~\ref{lem:poincare3-pi1-sum},
dostajemy
\[
 1=\pi_1(M)\cong
 \Gamma_1*\cdots*\Gamma_k*
 \underbrace{\Z*\cdots*\Z}_{\ell\ \text{razy}}.
\]
Każdy czynnik kanonicznie wstrzykuje się
w wolny iloczyn. Grupa po prawej może
więc być trywialna tylko wtedy, gdy
$\ell=0$ i każda $\Gamma_i=1$.
Wtedy każdy pozostały czynnik jest
$S^3/\{1\}=S^3$, a suma spójna
ze sferą $S^3$ nie zmienia dyfeomorfizmu
rozmaitości. Zatem $M\cong S^3$.
\end{proof}

Ostatni dowód jest pełnym \emph{wyprowadzeniem
wniosku topologicznego z wymienionych twierdzeń
zewnętrznych}. Nie jest pełnym dowodem
Twierdzenia~\ref{thm:poincare-trzy} od podstaw,
ponieważ twierdzenia o chirurgii i wygasaniu
przyjęliśmy bez dowodu. Jest to ta sama
różnica między użyciem a udowodnieniem
głębokiego wyniku, którą zaznaczaliśmy
w Rozdziale~\ref{ch:poincare-2-high}.

\begin{wniosek}[Topologiczna hipoteza Poincarégo w wymiarze trzy]
\label{cor:poincare-trzy-topologiczna}
Każda zamknięta, spójna i jednospójna topologiczna
$3$-rozmaitość jest homeomorficzna z $S^3$.
\end{wniosek}
\begin{proof}
Twierdzenie Moise'a zapewnia, że każda topologiczna
$3$-rozmaitość dopuszcza zgodną strukturę gładką;
jest to użyty tutaj wynik zewnętrzny.
Źródło: \href{https://msp.org/gtm/2015/19-1/gtm-v19-n1-p08-s.pdf}
{M. Aschenbrenner, S. Friedl i H. Wilton,
\emph{3-Manifold Groups}, twierdzenie 3.1},
z odwołaniem do oryginalnej pracy
\href{https://doi.org/10.2307/1969769}
{E.~E. Moise'a (1952)}. Po wyborze takiej struktury
Twierdzenie~\ref{thm:poincare-trzy} daje dyfeomorfizm
z $S^3$, a więc również homeomorfizm pierwotnej
rozmaitości z $S^3$.
\end{proof}

\section{Miejsce tego argumentu w książce}
\label{sec:poincare3-miejsce}

Twierdzenie Poincarégo w wymiarze trzy jest
wnioskiem o \emph{gładkim typie} rozmaitości.
W Rozdziale~\ref{ch:poincare-2-high} za punkt
wyjścia przyjmowaliśmy sferę homotopijną;
tutaj wystarcza jednospójność. Lemat
~\ref{lem:poincare3-pi3} pokazuje, że
zamknięta jednospójna $3$-rozmaitość ma
te same grupy homologii co $S^3$.
Właściwy krok od tych niezmienników do
dyfeomorfizmu wykonuje jednak analiza
potoku z chirurgią.

Słowo \emph{chirurgia} występuje w dwóch
powiązanych sensach. W rozdziałach
~\ref{ch:chirurgia-geometryczna-21}--
\ref{ch:ciag-chirurgii} badaliśmy topologiczne
operacje i przeszkody dla odwzorowań normalnych.
Tutaj chirurgia jest zmianą metrycznego potoku
w bardzo szczególnych zanurzonych szyjkach
$S^2\times I$. Lematy~\ref{lem:poincare3-sphere-cut}
i~\ref{lem:poincare3-pi1-sum} pozwalają
śledzić jej skutek topologiczny; otoczenia
kanoniczne i brak zapadania wskazują,
gdzie i kiedy wolno ją wykonać.

W szerszej teorii analiza długoczasowego
potoku z chirurgią prowadzi do geometryzacji
Thurstona. Dla samego twierdzenia Poincarégo
wystarcza gałąź argumentu kończąca się
wygasaniem w skończonym czasie. Dzięki temu
nie musimy tu klasyfikować geometrii
składowych, które przetrwałyby dla
$t\to\infty$.

\par\smallskip
{\footnotesize\noindent\emph{Źródła głównych wejść:}
\href{https://arxiv.org/abs/math/0211159}
{G.\ Perelman, \emph{The entropy formula for the Ricci
flow and its geometric applications} (2002)}:
entropia, niezapadanie i geometria granic;
\href{https://arxiv.org/abs/math/0303109}
{G.\ Perelman, \emph{Ricci flow with surgery on
three-manifolds} (2003)}: konstrukcja chirurgii;
\href{https://arxiv.org/abs/math/0307245}
{G.\ Perelman, \emph{Finite extinction time for
the solutions to the Ricci flow on certain
three-manifolds} (2003)}: wygasanie;
\href{https://www.claymath.org/library/monographs/cmim03.pdf}
{J.\ Morgan i G.\ Tian, \emph{Ricci Flow and the
Poincaré Conjecture}, Clay Mathematics Monographs 3
(2007)}: pełny szczegółowy dowód, zwłaszcza
rozdziały 15--19.\par}
